\documentclass[reqno]{amsart} 
\usepackage{amssymb,latexsym,amsmath,amscd,graphicx,setspace,amsthm,verbatim,enumitem}
\usepackage[margin = 3 cm]{geometry}

\def\upint{\mathchoice%
    {\mkern13mu\overline{\vphantom{\intop}\mkern7mu}\mkern-20mu}%
    {\mkern7mu\overline{\vphantom{\intop}\mkern7mu}\mkern-14mu}%
    {\mkern7mu\overline{\vphantom{\intop}\mkern7mu}\mkern-14mu}%
    {\mkern7mu\overline{\vphantom{\intop}\mkern7mu}\mkern-14mu}%
  \int}
\def\lowint{\mkern3mu\underline{\vphantom{\intop}\mkern7mu}\mkern-10mu\int}


\theoremstyle{plain}
\newtheorem{theorem}{Theorem}[section]
\newtheorem{proposition}{Proposition}
\newtheorem{corollary}{Corollary}
\newtheorem{lemma}{Lemma}
\newtheorem{conjecture}{Conjecture}
\newtheorem{question}{Question}
\newtheorem{problem}{Problem}
      
\theoremstyle{definition}
\newtheorem{definition}{Definition}

\newenvironment{solution1}{\paragraph{\emph{Solution $1$}.}}{\hfill$\square$}
\newenvironment{solution2}{\paragraph{\emph{Solution $2$}.}}{\hfill$\square$}
\newenvironment{solution3}{\paragraph{\emph{Solution $3$}.}}{\hfill$\square$}

\begin{document} 

\title[Homework 12]{Homework 12}

\date{\today} 
\maketitle 


In class, we reminded ourselves about the notion of lim sup and lim inf:  If $(a_{n})_{n=1}^{\infty}$ is a sequence such that $a_{n} \in [-\infty,\infty]$ for all $n \in \mathbb{N}$, then letting
$$b_{n} = \inf_{k \ge n }a_{k}  \text{ and } c_{n} = \sup_{k \ge n} a_{k}$$
one defines
$$\liminf_{n \to \infty}a_{n} = \lim_{n \to \infty}b_{n} \text{ and } \limsup_{n \to \infty}a_{n} = \lim_{n \to \infty}c_{n}, $$
and that makes sense since both sequences $(b_{n})_{n=1}^{\infty}$ and $(c_{n})_{n=1}^{\infty}$ are monotonic and hence converge in $[-\infty,\infty]$.

\begin{problem}
Show that
\begin{enumerate}
\item One always has
$$\liminf_{n \to \infty}a_{n} \le \limsup_{n \to \infty}a_{n}. $$
\item One has that $(a_{n})_{n=1}^{\infty}$ converges if and only if 
$$\liminf_{n \to \infty}a_{n} = \limsup_{n \to \infty}a_{n} $$
in which case the limit of $(a_{n})_{n=1}^{\infty}$ is precisely this common value.
\end{enumerate}
\end{problem}
\begin{solution1}

\end{solution1}

\begin{problem}
Let $(X,\mathcal{A},\mu)$ be a measure space and let
$$f:X \rightarrow [-\infty,\infty] $$
be measurable for which its Lebesgue integral is defined.  Show that if $A,B \in \mathcal{A}$ are such that
$$X = A \cup B \text{ and } \mu(A \cap B) = 0, $$
then both
$$\int_{A}f \text{ and } \int_{B}f $$
are defined and 
$$\int_{X}f = \int_{A}f + \int_{B}f. $$
\end{problem}
\begin{solution1}

\end{solution1}


\begin{problem}
In class, we went over the following example of a pointwise converging sequence of functions.  The sequence was $f_{n}:[0,1] \rightarrow [0,\infty]$ defined via
\begin{equation*}
x \mapsto f_{n}(x) = 
\begin{cases}
n,& \text{ if } 0 \le x \le \frac{1}{n} ; \\
1-x, & \text{ if } \frac{1}{n} \le x \le 1.
\end{cases}
\end{equation*}
In homework 11, you studied this sequence of functions carefully.  Does it contradict the dominated convergence theorem??  Explain carefully.
\end{problem}
\begin{solution1}

\end{solution1}

\noindent\rule{\linewidth}{0.4pt}

The rest of this problem set will consist in realizing that the theory of Lebesgue integration contains as a particular case the theory of absolutely convergent series.  We start with the following general definition:

\begin{definition}
Let $(X,\mathcal{A},\mu)$ be a measure space and let 
$$f: X \rightarrow \mathbb{R} $$
be a measurable function.  One says that $f$ is {\bfseries Lebesgue integrable} if 
$$\int_{X} |f| < \infty. $$
The collection of all Lebesgue integrable functions is denoted by $\mathcal{L}^{1}(X,\mu)$.
\end{definition}
Be sure to realize that saying that $f$ is Lebesgue integrable is {\bfseries stronger} than just assuming that its Lebesgue integral is defined.

\begin{problem}
Put your linear algebra hat on and show that $\mathcal{L}^{1}(X,\mu)$ is a vector subspace of $\mathbb{R}^{X}$ so that $\mathcal{L}^{1}(X,\mu)$ is in particular a vector space over $\mathbb{R}$ itself.  (The set $\mathbb{R}^{X}$ denotes the vector space of real-valued functions on $X$.)
\end{problem}
\begin{solution1}

\end{solution1}

Recall that in class we pointed out that $(\mathbb{N},\mathcal{P}(\mathbb{N}),\mu)$, where $\mu:\mathcal{P}(\mathbb{N}) \rightarrow [0,\infty]$ is the {\bfseries counting measure} is a measure space.

\begin{problem}
\hfill
\begin{enumerate}
\item Explain why every function $f:\mathbb{N} \rightarrow [-\infty,\infty]$ is measurable.
\item What subsets of $\mathbb{N}$ have measure zero??
\end{enumerate}
\end{problem}
\begin{solution1}

\end{solution1}

\begin{problem}
Let $a_{1},\ldots,a_{n} \in [0,\infty)$ and define $f:\mathbb{N} \rightarrow [0,\infty)$ via
$$i \mapsto f(i) = \begin{cases} a_{i}, &\text{ if } 1 \le i \le n;\\ 0, &\text{ if } i > n.  \end{cases} $$
Show that
$$\int_{\mathbb{N}}f = \sum_{i=1}^{n}a_{i}. $$
\end{problem}
\begin{solution1}

\end{solution1}

\begin{problem}
Let $(a_{i})_{i=1}^{\infty}$ be a sequence such that $a_{i} \in [0,\infty)$ for all $i \in \mathbb{N}$ and define $f_{n}:\mathbb{N} \rightarrow [0,\infty)$ via
$$i \mapsto f_{n}(i) = \begin{cases} a_{i}, & \text{ if } 1 \le i \le n;\\0, & \text{ if } i < n. \end{cases} $$
\begin{enumerate}
\item Show that the sequence $(f_{n})_{n=1}^{\infty}$ converges pointwise to a function $f$ and determine $f$.
\item Show that the sequence $(f_{n})_{n=1}^{\infty}$ is increasing and apply the monotone convergence theorem to deduce that
$$\int_{\mathbb{N}} f = \sum_{i=1}^{\infty}a_{i}. $$
\end{enumerate}
\end{problem}

Given a sequence $(a_{i})_{i=1}^{\infty}$ of real numbers, one defines
$$a_{i}^{+} = \begin{cases} a_{i}, \text{ if } a_{i} \ge 0; \\ 0, \text{ if } a_{i} < 0, \end{cases} \text{ and also } a_{i}^{-} = \begin{cases} 0, \text{ if } a_{i} \ge 0; \\ -a_{i}, \text{ if } a_{i} < 0. \end{cases} $$

\begin{problem}
Let $(a_{i})_{i=1}^{\infty}$ be a sequence such that $a_{i} \in \mathbb{R}$ for all $i \in \mathbb{R}$ and define $f:\mathbb{N} \rightarrow \mathbb{R}$ via 
$$i \mapsto f(i) = a_{i}. $$
\begin{enumerate}
\item Show that
$$\int_{\mathbb{N}}f^{+} = \sum_{i=1}^{\infty}a_{i}^{+} \text{ and } \int_{\mathbb{N}}f^{-} = \sum_{i=1}^{\infty}a_{i}^{-}. $$
\item Show that $f \in \mathcal{L}^{1}(\mathbb{N},\mu)$ if and only if the series 
$$\sum_{i=1}^{\infty}a_{i} $$
converges {\bfseries absolutely}.
\item Show that if $f \in \mathcal{L}^{1}(\mathbb{N},\mu)$, then
$$\int_{\mathbb{N}}f = \sum_{i=1}^{\infty}a_{i}. $$
\item Assume that $a_{i} = (-1)^{i+1}/i$.  Does the corresponding series $\sum a_{i}$ converge?  Does one have $f \in \mathcal{L}^{1}(\mathbb{N},\mu)$?
\end{enumerate}

\end{problem}

For the particular case of the measure space $(\mathbb{N},\mathcal{P}(\mathbb{N}),\mu)$, where $\mu$ is the counting measure, one often denotes $\mathcal{L}^{1}(\mathbb{N},\mu)$ simply by $\ell^{1}$.  Thus by the last problem, one has:
$$\ell^{1} = \left\{ (a_{i})_{i=1}^{\infty}: a_{i}\in \mathbb{R} \text{ for all } i \in \mathbb{R} \text{ and } \sum_{i=1}^{\infty}|a_{i}|< \infty \right\}. $$

Start thinking about what the dominated convergence theorem would say for series...I'll keep this for the next homework set.

\end{document}









